% Section 2 -- drafted by Claude (2026-10-10), revised after a fact check (2026-10-11); V5 rules added 2026-10-21.
\section{Task, data and local evaluation}\label{sec:data}

The predictor outputs the probability that a subject succeeds on an item. A subject is an AI system defined by its model, evaluation harness, harness version and reasoning effort; the items come from benchmarks in the competition's private test pool \citep{truong2026paiec}. Codabench \citep{xu2022codabench} scores a draw of subject--benchmark pairs. Each pair's items are split in half: an acquisition stream that supplies the labels and an evaluation set that is scored after budgets of $b \in \{0, 1, 3, 7, 15, 31\}$ acquired labels. With $B_b$ the pair's Brier score \citep{brier1950} at budget $b$, the pair's Brier ALC (area under the learning curve) is
\begin{equation}
  \mathrm{ALC} = 0.1\,B_0 + 0.2\,(B_1 + B_3 + B_7 + B_{15}) + 0.1\,B_{31},
  \label{eq:alc}
\end{equation}
and the score of a submission is the mean over pairs. Lower is better. At each budget the predictor receives the labels acquired so far for every pair in the draw (none at budget 0), so when a draw puts several subjects on one benchmark, it can use the other subjects' labels. None of my uploads included a labeling policy, and the platform acquired labels at random. Each upload gets formative feedback: its score on one draw of pairs (8 or 9 for me), drawn afresh each time, with per-budget Brier tables.

\paragraph{Data.}
I trained and tested on the public measurement-db \citep{measurementdb2026} at revision \texttt{bc8204d8}; the only other training data are the LLM difficulty ratings and subject guesses described in Section~\ref{sec:method}. Because the competition scores binary outcomes, I kept an item only if its declared response scale is $\{0,1\}$. This drops MMDocRAG \citep{dong2025mmdocrag}, graded on a 26-level scale, and 271 MathArena items graded on an interval scale (Appendix~\ref{app:harness}). What is left is 10,361 items, 229 subjects and 225,764 responses on five benchmarks: MathArena \citep{balunovic2025matharena}, Multi-SWE-bench \citep{zan2025multiswe}, REAL \citep{garg2025real}, ResearchCodeBench \citep{hua2025researchcodebench} and SWE-rebench \citep{badertdinov2025swerebench} (Table~\ref{tab:data}).

\begin{table}[h]
  \caption{Public training data after the binary-scale filter. Responses include repeated trials; mean label is the share equal to 1. Pairs: pairs with 80 to 1,000 items (default setting, per seed). Rated items: items with an LLM rating (1,379 items share 1,049 ratings; on MathArena, 696 items cover 366 problems). Rated pairs: pairs with at least 40 rated items (rated setting, per seed). The total row counts distinct subjects.}
  \label{tab:data}
  \centering
  \small
  \setlength{\tabcolsep}{4pt}
  \begin{tabular}{lrrrrrrr}
\toprule
Benchmark & Items & Responses & Subjects & Mean label & Pairs & Rated items & Rated pairs \\
\midrule
MathArena & 1,484 & 89,482 & 95 & 0.67 & 81 & 696 & 65 \\
Multi-SWE-bench & 2,126 & 57,975 & 82 & 0.15 & 54 & 400 & 70 \\
REAL & 233 & 4,449 & 33 & 0.36 & 26 & 233 & 32 \\
ResearchCodeBench & 212 & 6,784 & 31 & 0.36 & 31 & 50 & 31 \\
SWE-rebench & 6,306 & 67,074 & 1 & 0.48 & 0 & -- & -- \\
\midrule
Total & 10,361 & 225,764 & 229 & & 192 & 1,379 & 198 \\
\bottomrule
\end{tabular}

\end{table}

\paragraph{Local evaluation.}
To compare predictors offline I wrote a harness that simulates the hosted protocol on the public data. Each fold holds out one benchmark: the predictor is fit on the other four and scored on pairs from the held-out one. Each subject with enough items on the held-out benchmark forms a pair, whose items are shuffled with a seeded generator and split in half. Labels are acquired in rounds (Appendix~\ref{app:harness}) with the default random rule of the competition's starting kit, and at each budget the predictor is scored on every recorded response of the evaluation items. Recomputing each pair's ALC from the per-budget Brier tables of my Codabench feedback reproduces the reported value to within $10^{-6}$ for all 68 pair results I kept. Appendix~\ref{app:harness} gives the details and the check's limits.

I compare predictors in three settings. The \emph{default} setting follows the rules' minimum of 80 items per pair and caps pairs at 1,000 items, the platform's limit for a whole draw; it gives 192 pairs per seed, with 265 items on average. The \emph{hidden-sized} setting imitates the smaller pairs of a Codabench draw, where the 8 or 9 pairs of each draw I saw shared at most 1,000 subject--item combinations: it samples at most 400 items of the held-out benchmark and keeps pairs with at least 60, about 157 pairs per seed with 153 items on average. For the \emph{rated} setting, which tests the LLM inputs, I kept only items with an LLM rating and pairs with at least 40 of them (198 pairs per seed). I lowered the minimum below the rules' 80 in the last two settings to keep enough pairs. In any setting a pair can also be given the labels of $K$ partner pairs on the same benchmark, and I give $K$ with each result.

All runs use seeds 0, 1 and 2, which set the split, the stream order and the partners but not the fitted predictor. A comparison reports $\Delta = \mathrm{ALC}(\text{candidate}) - \mathrm{ALC}(\text{reference})$, averaged over all seeds' pair results, matched on seed and pair, so $\Delta < 0$ favors the candidate. Its standard error treats each subject--benchmark pair as one cluster across the seeds (CR0, Appendix~\ref{app:harness}). Appendix~\ref{app:rules} lists each test's rule, written before its run, and the choices made without a rule.
